\section{Historical appendix: Scriptural Date and Location}
\label{historical-appendix}
% Generated. Do not edit.
% Deterministic projection of research/chronology.toml.
% schema: 5
% document: liturgy/roman-rite/1962/propers/temporal/66-resumed-fifth-after-epiphany
% generated-by: tools/tpt proper-chronology annotations
\newcommand{\chronologyannotationclaim}[6]{#6}
\newcommand{\chronologyannotationcomparisonclaim}[8]{#8}
\newcommand{\chronologyannotationreach}[3]{}
\newcommand{\chronologyannotationgroup}[3]{#3}
\newcommand{\chronologyannotationcomparisongroup}[4]{#4}
\newcommand{\chronologyannotation}[1]{%
  \ifcsname triptychchronologyannotation@#1\endcsname
    \csname triptychchronologyannotation@#1\endcsname
  \else
    \PackageError{triptych}{No chronology annotation for #1}{}%
  \fi}
% comparison-dependency: src/gpt/liturgy/roman-rite/1962/propers/temporal/66-resumed-fifth-after-epiphany/research/chronology-profile-comparisons.toml sha256=0556d3276a9de8069099044e7830e0ad093924429fc6375de60c4481511e2788
% introit: Introit
\expandafter\def\csname triptychchronologyannotation@introit\endcsname{%
  \chronologyannotationgroup{historical-setting}{research-pending}{\textbf{Historical setting} -- Occasion date unresolved.}%
  \space \chronologyannotationgroup{composition}{nonuniform}{\textbf{Composition} -- No single date applies across every cited locus.}%
}
% epistle: Epistle
\expandafter\def\csname triptychchronologyannotation@epistle\endcsname{%
  \chronologyannotationgroup{composition}{disputed}{\textbf{Composition} -- disputed: \chronologyannotationreach{Col.3.12}{Col}{true}\chronologyannotationreach{Col.3.13}{Col}{true}\chronologyannotationreach{Col.3.14}{Col}{true}\chronologyannotationreach{Col.3.15}{Col}{true}\chronologyannotationreach{Col.3.16}{Col}{true}\chronologyannotationreach{Col.3.17}{Col}{true}\chronologyannotationclaim{composition.epistle-to-the-colossians}{composition}{catholic-traditional-v1}{disputed}{(Philemon; Colossians; Ephesians; Philippians), 61}{A.D. 61}; \chronologyannotationreach{Col.3.12}{Col}{true}\chronologyannotationreach{Col.3.13}{Col}{true}\chronologyannotationreach{Col.3.14}{Col}{true}\chronologyannotationreach{Col.3.15}{Col}{true}\chronologyannotationreach{Col.3.16}{Col}{true}\chronologyannotationreach{Col.3.17}{Col}{true}\chronologyannotationclaim{composition.epistle-to-the-colossians}{composition}{catholic-traditional-v1}{disputed}{some time between A.D. 61 and 63}{A.D. 61--63}.}%
}
% gradual: Gradual
\expandafter\def\csname triptychchronologyannotation@gradual\endcsname{%
  \chronologyannotationgroup{historical-setting}{research-pending}{\textbf{Historical setting} -- Occasion date unresolved.}%
  \space \chronologyannotationgroup{composition}{preferred}{\textbf{Composition}: \chronologyannotationreach{Ps.43.8}{Ps.41 Ps.42 Ps.43 Ps.44 Ps.45 Ps.46 Ps.47 Ps.48 Ps.83 Ps.84 Ps.86 Ps.87}{true}\chronologyannotationreach{Ps.43.9}{Ps.41 Ps.42 Ps.43 Ps.44 Ps.45 Ps.46 Ps.47 Ps.48 Ps.83 Ps.84 Ps.86 Ps.87}{true}\chronologyannotationclaim{composition.psalms-of-the-sons-of-korah}{composition}{catholic-traditional-v1}{preferred}{between the days of Isaias and the return from exile}{Between the days of Isaias and the return from exile}.}%
}
% alleluia: Alleluia
\expandafter\def\csname triptychchronologyannotation@alleluia\endcsname{%
  \chronologyannotationgroup{historical-setting}{research-pending}{\textbf{Historical setting} -- Occasion date unresolved.}%
  \space \chronologyannotationgroup{composition}{preferred}{\textbf{Composition}: \chronologyannotationreach{Ps.129.1}{Ps}{true}\chronologyannotationreach{Ps.129.2}{Ps}{true}\chronologyannotationclaim{critical.psalms.latest-composition-boundary}{composition}{catholic-critical-v1}{preferred}{before the Maccabean period, around 165 B.C.; no individual psalm can be dated securely}{Before c. 165 B.C.}}%
}
% gospel: Gospel
\expandafter\def\csname triptychchronologyannotation@gospel\endcsname{%
  \chronologyannotationgroup{narrated-event}{research-pending}{\textbf{Event} -- Narrated event date unresolved.}%
  \space \chronologyannotationgroup{composition}{disputed}{\textbf{Composition} -- disputed: \chronologyannotationreach{Matt.13.24}{Matt}{true}\chronologyannotationreach{Matt.13.25}{Matt}{true}\chronologyannotationreach{Matt.13.26}{Matt}{true}\chronologyannotationreach{Matt.13.27}{Matt}{true}\chronologyannotationreach{Matt.13.28}{Matt}{true}\chronologyannotationreach{Matt.13.29}{Matt}{true}\chronologyannotationreach{Matt.13.30}{Matt}{true}\chronologyannotationclaim{composition.gospel-of-matthew}{composition}{catholic-traditional-v1}{disputed}{about A.D. 38-45}{c. A.D. 38--45}; \chronologyannotationreach{Matt.13.24}{Matt}{true}\chronologyannotationreach{Matt.13.25}{Matt}{true}\chronologyannotationreach{Matt.13.26}{Matt}{true}\chronologyannotationreach{Matt.13.27}{Matt}{true}\chronologyannotationreach{Matt.13.28}{Matt}{true}\chronologyannotationreach{Matt.13.29}{Matt}{true}\chronologyannotationreach{Matt.13.30}{Matt}{true}\chronologyannotationclaim{composition.gospel-of-matthew}{composition}{catholic-traditional-v1}{disputed}{about the year 40-42}{c. A.D. 40--42}; \chronologyannotationreach{Matt.13.24}{Matt}{true}\chronologyannotationreach{Matt.13.25}{Matt}{true}\chronologyannotationreach{Matt.13.26}{Matt}{true}\chronologyannotationreach{Matt.13.27}{Matt}{true}\chronologyannotationreach{Matt.13.28}{Matt}{true}\chronologyannotationreach{Matt.13.29}{Matt}{true}\chronologyannotationreach{Matt.13.30}{Matt}{true}\chronologyannotationclaim{composition.gospel-of-matthew}{composition}{catholic-traditional-v1}{disputed}{the years 40-45}{A.D. 40--45}; \chronologyannotationreach{Matt.13.24}{Matt}{true}\chronologyannotationreach{Matt.13.25}{Matt}{true}\chronologyannotationreach{Matt.13.26}{Matt}{true}\chronologyannotationreach{Matt.13.27}{Matt}{true}\chronologyannotationreach{Matt.13.28}{Matt}{true}\chronologyannotationreach{Matt.13.29}{Matt}{true}\chronologyannotationreach{Matt.13.30}{Matt}{true}\chronologyannotationclaim{composition.gospel-of-matthew}{composition}{catholic-traditional-v1}{disputed}{about the year 60-68}{c. A.D. 60--68}; \chronologyannotationreach{Matt.13.24}{Matt}{true}\chronologyannotationreach{Matt.13.25}{Matt}{true}\chronologyannotationreach{Matt.13.26}{Matt}{true}\chronologyannotationreach{Matt.13.27}{Matt}{true}\chronologyannotationreach{Matt.13.28}{Matt}{true}\chronologyannotationreach{Matt.13.29}{Matt}{true}\chronologyannotationreach{Matt.13.30}{Matt}{true}\chronologyannotationclaim{composition.gospel-of-matthew}{composition}{catholic-traditional-v1}{disputed}{about the years 64-67}{c. A.D. 64--67}; \chronologyannotationreach{Matt.13.24}{Matt}{true}\chronologyannotationreach{Matt.13.25}{Matt}{true}\chronologyannotationreach{Matt.13.26}{Matt}{true}\chronologyannotationreach{Matt.13.27}{Matt}{true}\chronologyannotationreach{Matt.13.28}{Matt}{true}\chronologyannotationreach{Matt.13.29}{Matt}{true}\chronologyannotationreach{Matt.13.30}{Matt}{true}\chronologyannotationclaim{composition.gospel-of-matthew}{composition}{catholic-traditional-v1}{disputed}{about the year 50}{c. A.D. 50}.}%
  \space \chronologyannotationcomparisongroup{catholic-critical-v1}{composition}{composition-only}{\textbf{Composition (Catholic critical chronology, for comparison)}: \chronologyannotationreach{Matt.13.24}{Matt}{true}\chronologyannotationreach{Matt.13.25}{Matt}{true}\chronologyannotationreach{Matt.13.26}{Matt}{true}\chronologyannotationreach{Matt.13.27}{Matt}{true}\chronologyannotationreach{Matt.13.28}{Matt}{true}\chronologyannotationreach{Matt.13.29}{Matt}{true}\chronologyannotationreach{Matt.13.30}{Matt}{true}\chronologyannotationcomparisonclaim{critical.gospel-of-matthew}{composition}{catholic-critical-v1}{catholic-critical-v1}{preferred}{passage.united-states-conference-of-catholic-bishops.new-american-bible-revised-edition.english-usccb-web-2026-09-21.matthew-introduction}{post-A.D. 70 date}{The Greek Gospel of Matthew in the NABRE introduction, post-A.D. 70 date}.}%
}
% offertory: Offertory
\expandafter\def\csname triptychchronologyannotation@offertory\endcsname{%
  \chronologyannotationgroup{historical-setting}{research-pending}{\textbf{Historical setting} -- Occasion date unresolved.}%
  \space \chronologyannotationgroup{composition}{preferred}{\textbf{Composition}: \chronologyannotationreach{Ps.129.1}{Ps}{true}\chronologyannotationreach{Ps.129.2}{Ps}{true}\chronologyannotationclaim{critical.psalms.latest-composition-boundary}{composition}{catholic-critical-v1}{preferred}{before the Maccabean period, around 165 B.C.; no individual psalm can be dated securely}{Before c. 165 B.C.}}%
}
% communion: Communion
\expandafter\def\csname triptychchronologyannotation@communion\endcsname{%
  \chronologyannotationgroup{composition}{disputed}{\textbf{Composition} -- disputed: \chronologyannotationreach{Mark.11.24}{Mark}{true}\chronologyannotationclaim{composition.gospel-of-mark}{composition}{catholic-traditional-v1}{disputed}{the tenth or twelfth year after the Ascension (A.D. 38-40)}{A.D. 38--40}; \chronologyannotationreach{Mark.11.24}{Mark}{true}\chronologyannotationclaim{composition.gospel-of-mark}{composition}{catholic-traditional-v1}{disputed}{A.D. 40}{A.D. 40}; \chronologyannotationreach{Mark.11.24}{Mark}{true}\chronologyannotationclaim{composition.gospel-of-mark}{composition}{catholic-traditional-v1}{disputed}{the third year of Claudius (A.D. 43)}{A.D. 43}; \chronologyannotationreach{Mark.11.24}{Mark}{true}\chronologyannotationclaim{composition.gospel-of-mark}{composition}{catholic-traditional-v1}{disputed}{somewhere between A.D. 50 and 67}{A.D. 50--67}; \chronologyannotationreach{Mark.11.24}{Mark}{true}\chronologyannotationclaim{composition.gospel-of-mark}{composition}{catholic-traditional-v1}{disputed}{between 60 and 70}{A.D. 60--70}; \chronologyannotationreach{Mark.11.24}{Mark}{true}\chronologyannotationclaim{composition.gospel-of-mark}{composition}{catholic-traditional-v1}{disputed}{after Matthew, and not as late as the destruction of Jerusalem}{After Matthew, and not as late as the destruction of Jerusalem}.}%
  \space \chronologyannotationcomparisongroup{catholic-critical-v1}{composition}{composition-only}{\textbf{Composition (Catholic critical chronology, for comparison)}: \chronologyannotationreach{Mark.11.24}{Mark.1 Mark.2 Mark.3 Mark.4 Mark.5 Mark.6 Mark.7 Mark.8 Mark.9 Mark.10 Mark.11 Mark.12 Mark.13 Mark.14 Mark.15 Mark.16.1-8}{true}\chronologyannotationcomparisonclaim{critical.gospel-of-mark}{composition}{catholic-critical-v1}{catholic-critical-v1}{preferred}{artifact.united-states-conference-of-catholic-bishops.new-american-bible-revised-edition.english-usccb-web-2026-11-02.mark-introduction-web-text-272bcce7}{perhaps shortly after the year 70}{The Gospel of Mark under the NABRE's tentative modern dating proposal, perhaps shortly after A.D. 70}.}%
}
% presentation: orientation-v2
\newcommand{\chronologyorientationrow}[2]{%
  \ifcsname triptychchronologyorientation@#1\endcsname
    \csname triptychchronologyorientation@#1\endcsname{#2}%
  \else
    \PackageError{triptych}{No chronology orientation row for #1}{}%
  \fi}
\expandafter\def\csname triptychchronologyorientation@row-4497fc13bd16e15c\endcsname#1{%
Introit: Ps 84:2 (Vulgate numbering; Hebrew 85); Gradual: Psalm 43:8--9 (Vulgate numbering; Hebrew 44); Alleluia: Psalm 129:1--2 (Vulgate numbering; Hebrew 130); Offertory: Psalm 129:1--2 (Vulgate numbering; Hebrew 130) & Historical setting. #1 & Occasion date unresolved\\[.25em]%
}
\expandafter\def\csname triptychchronologyorientation@row-d1f3edd1ff63c2d2\endcsname#1{%
Introit: Ps 84:2 (Vulgate numbering; Hebrew 85); Gradual: Psalm 43:8--9 (Vulgate numbering; Hebrew 44) & Composition: The Psalms of the Sons of Korah. #1 & \chronologyannotationreach{Ps.43.8}{Ps.41 Ps.42 Ps.43 Ps.44 Ps.45 Ps.46 Ps.47 Ps.48 Ps.83 Ps.84 Ps.86 Ps.87}{true}\chronologyannotationreach{Ps.43.9}{Ps.41 Ps.42 Ps.43 Ps.44 Ps.45 Ps.46 Ps.47 Ps.48 Ps.83 Ps.84 Ps.86 Ps.87}{true}\chronologyannotationreach{Ps.84.2}{Ps.41 Ps.42 Ps.43 Ps.44 Ps.45 Ps.46 Ps.47 Ps.48 Ps.83 Ps.84 Ps.86 Ps.87}{true}\chronologyannotationclaim{composition.psalms-of-the-sons-of-korah}{composition}{catholic-traditional-v1}{preferred}{between the days of Isaias and the return from exile}{Between the days of Isaias and the return from exile}\\[.25em]%
}
\expandafter\def\csname triptychchronologyorientation@row-da4404f3dbb16b26\endcsname#1{%
Alleluia: Psalm 129:1--2 (Vulgate numbering; Hebrew 130); Offertory: Psalm 129:1--2 (Vulgate numbering; Hebrew 130) & Composition: Psalms. #1 & \chronologyannotationreach{Ps.129.1}{Ps}{true}\chronologyannotationreach{Ps.129.2}{Ps}{true}\chronologyannotationclaim{critical.psalms.latest-composition-boundary}{composition}{catholic-critical-v1}{preferred}{before the Maccabean period, around 165 B.C.; no individual psalm can be dated securely}{Before c. 165 B.C.}\\[.25em]%
}
\expandafter\def\csname triptychchronologyorientation@row-84a964a3effacafd\endcsname#1{%
Introit: Jer 29:11--12,14 & Composition: Jeremias's letter to the exiles after Jechonias's deportation. #1 & \chronologyannotationreach{Jer.29.11}{Jer.29.1-23}{false}\chronologyannotationreach{Jer.29.12}{Jer.29.1-23}{false}\chronologyannotationreach{Jer.29.14}{Jer.29.1-23}{false}\chronologyannotationclaim{composition.jeremias-letter-to-the-exiles}{composition}{catholic-traditional-v1}{preferred}{After that Jechonias the king ... were departed out of Jerusalem}{After Jechonias and the deportees had departed from Jerusalem}\\[.25em]%
}
\expandafter\def\csname triptychchronologyorientation@row-cc5afc871115711e\endcsname#1{%
Introit: Jer 29:11--12,14 & Final formation: The received Book of Jeremias after enlargement. #1 & \chronologyannotationreach{Jer.29.11}{Jer}{true}\chronologyannotationreach{Jer.29.12}{Jer}{true}\chronologyannotationreach{Jer.29.14}{Jer}{true}\chronologyannotationclaim{composition.book-of-jeremias}{final-formation}{catholic-traditional-v1}{preferred}{the first transcript was enlarged by various insertions and additions}{After the first authentic recension}\\[.25em]%
}
\expandafter\def\csname triptychchronologyorientation@row-a1420e3e43d2b570\endcsname#1{%
Gospel: Matthew 13:24--30 & Event. #1 & Date unresolved\\[.25em]%
}
\expandafter\def\csname triptychchronologyorientation@row-68d3fd83119cd2be\endcsname#1{%
Gospel: Matthew 13:24--30 & Composition: The Gospel of St Matthew. #1 & Disputed: \chronologyannotationreach{Matt.13.24}{Matt}{true}\chronologyannotationreach{Matt.13.25}{Matt}{true}\chronologyannotationreach{Matt.13.26}{Matt}{true}\chronologyannotationreach{Matt.13.27}{Matt}{true}\chronologyannotationreach{Matt.13.28}{Matt}{true}\chronologyannotationreach{Matt.13.29}{Matt}{true}\chronologyannotationreach{Matt.13.30}{Matt}{true}\chronologyannotationclaim{composition.gospel-of-matthew}{composition}{catholic-traditional-v1}{disputed}{about A.D. 38-45}{c. A.D. 38--45}; \chronologyannotationreach{Matt.13.24}{Matt}{true}\chronologyannotationreach{Matt.13.25}{Matt}{true}\chronologyannotationreach{Matt.13.26}{Matt}{true}\chronologyannotationreach{Matt.13.27}{Matt}{true}\chronologyannotationreach{Matt.13.28}{Matt}{true}\chronologyannotationreach{Matt.13.29}{Matt}{true}\chronologyannotationreach{Matt.13.30}{Matt}{true}\chronologyannotationclaim{composition.gospel-of-matthew}{composition}{catholic-traditional-v1}{disputed}{about the year 40-42}{c. A.D. 40--42}; \chronologyannotationreach{Matt.13.24}{Matt}{true}\chronologyannotationreach{Matt.13.25}{Matt}{true}\chronologyannotationreach{Matt.13.26}{Matt}{true}\chronologyannotationreach{Matt.13.27}{Matt}{true}\chronologyannotationreach{Matt.13.28}{Matt}{true}\chronologyannotationreach{Matt.13.29}{Matt}{true}\chronologyannotationreach{Matt.13.30}{Matt}{true}\chronologyannotationclaim{composition.gospel-of-matthew}{composition}{catholic-traditional-v1}{disputed}{the years 40-45}{A.D. 40--45}; \chronologyannotationreach{Matt.13.24}{Matt}{true}\chronologyannotationreach{Matt.13.25}{Matt}{true}\chronologyannotationreach{Matt.13.26}{Matt}{true}\chronologyannotationreach{Matt.13.27}{Matt}{true}\chronologyannotationreach{Matt.13.28}{Matt}{true}\chronologyannotationreach{Matt.13.29}{Matt}{true}\chronologyannotationreach{Matt.13.30}{Matt}{true}\chronologyannotationclaim{composition.gospel-of-matthew}{composition}{catholic-traditional-v1}{disputed}{about the year 60-68}{c. A.D. 60--68}; \chronologyannotationreach{Matt.13.24}{Matt}{true}\chronologyannotationreach{Matt.13.25}{Matt}{true}\chronologyannotationreach{Matt.13.26}{Matt}{true}\chronologyannotationreach{Matt.13.27}{Matt}{true}\chronologyannotationreach{Matt.13.28}{Matt}{true}\chronologyannotationreach{Matt.13.29}{Matt}{true}\chronologyannotationreach{Matt.13.30}{Matt}{true}\chronologyannotationclaim{composition.gospel-of-matthew}{composition}{catholic-traditional-v1}{disputed}{about the years 64-67}{c. A.D. 64--67}; \chronologyannotationreach{Matt.13.24}{Matt}{true}\chronologyannotationreach{Matt.13.25}{Matt}{true}\chronologyannotationreach{Matt.13.26}{Matt}{true}\chronologyannotationreach{Matt.13.27}{Matt}{true}\chronologyannotationreach{Matt.13.28}{Matt}{true}\chronologyannotationreach{Matt.13.29}{Matt}{true}\chronologyannotationreach{Matt.13.30}{Matt}{true}\chronologyannotationclaim{composition.gospel-of-matthew}{composition}{catholic-traditional-v1}{disputed}{about the year 50}{c. A.D. 50}\\[.25em]%
}
\expandafter\def\csname triptychchronologyorientation@row-941543664efaa9ec\endcsname#1{%
Gospel: Matthew 13:24--30 & Composition (Catholic critical chronology, for comparison): The Greek Gospel of Matthew in the NABRE introduction. #1 & \chronologyannotationreach{Matt.13.24}{Matt}{true}\chronologyannotationreach{Matt.13.25}{Matt}{true}\chronologyannotationreach{Matt.13.26}{Matt}{true}\chronologyannotationreach{Matt.13.27}{Matt}{true}\chronologyannotationreach{Matt.13.28}{Matt}{true}\chronologyannotationreach{Matt.13.29}{Matt}{true}\chronologyannotationreach{Matt.13.30}{Matt}{true}\chronologyannotationcomparisonclaim{critical.gospel-of-matthew}{composition}{catholic-critical-v1}{catholic-critical-v1}{preferred}{passage.united-states-conference-of-catholic-bishops.new-american-bible-revised-edition.english-usccb-web-2026-09-21.matthew-introduction}{post-A.D. 70 date}{Post-A.D. 70 date}\\[.25em]%
}
\expandafter\def\csname triptychchronologyorientation@row-8b7f20d842814860\endcsname#1{%
Communion: Mark 11:24 & Composition: The Gospel of St Mark. #1 & Disputed: \chronologyannotationreach{Mark.11.24}{Mark}{true}\chronologyannotationclaim{composition.gospel-of-mark}{composition}{catholic-traditional-v1}{disputed}{the tenth or twelfth year after the Ascension (A.D. 38-40)}{A.D. 38--40}; \chronologyannotationreach{Mark.11.24}{Mark}{true}\chronologyannotationclaim{composition.gospel-of-mark}{composition}{catholic-traditional-v1}{disputed}{A.D. 40}{A.D. 40}; \chronologyannotationreach{Mark.11.24}{Mark}{true}\chronologyannotationclaim{composition.gospel-of-mark}{composition}{catholic-traditional-v1}{disputed}{the third year of Claudius (A.D. 43)}{A.D. 43}; \chronologyannotationreach{Mark.11.24}{Mark}{true}\chronologyannotationclaim{composition.gospel-of-mark}{composition}{catholic-traditional-v1}{disputed}{somewhere between A.D. 50 and 67}{A.D. 50--67}; \chronologyannotationreach{Mark.11.24}{Mark}{true}\chronologyannotationclaim{composition.gospel-of-mark}{composition}{catholic-traditional-v1}{disputed}{between 60 and 70}{A.D. 60--70}; \chronologyannotationreach{Mark.11.24}{Mark}{true}\chronologyannotationclaim{composition.gospel-of-mark}{composition}{catholic-traditional-v1}{disputed}{after Matthew, and not as late as the destruction of Jerusalem}{After Matthew, and not as late as the destruction of Jerusalem}\\[.25em]%
}
\expandafter\def\csname triptychchronologyorientation@row-f35007a3f1f636d7\endcsname#1{%
Communion: Mark 11:24 & Composition (Catholic critical chronology, for comparison): The Gospel of Mark under the NABRE's tentative modern dating proposal. #1 & \chronologyannotationreach{Mark.11.24}{Mark.1 Mark.2 Mark.3 Mark.4 Mark.5 Mark.6 Mark.7 Mark.8 Mark.9 Mark.10 Mark.11 Mark.12 Mark.13 Mark.14 Mark.15 Mark.16.1-8}{true}\chronologyannotationcomparisonclaim{critical.gospel-of-mark}{composition}{catholic-critical-v1}{catholic-critical-v1}{preferred}{artifact.united-states-conference-of-catholic-bishops.new-american-bible-revised-edition.english-usccb-web-2026-11-02.mark-introduction-web-text-272bcce7}{perhaps shortly after the year 70}{Perhaps shortly after A.D. 70}\\[.25em]%
}
\expandafter\def\csname triptychchronologyorientation@row-6be86a5d4d4ef64b\endcsname#1{%
Epistle: Colossians 3:12--17 & Composition: The Epistle to the Colossians. #1 & Disputed: \chronologyannotationreach{Col.3.12}{Col}{true}\chronologyannotationreach{Col.3.13}{Col}{true}\chronologyannotationreach{Col.3.14}{Col}{true}\chronologyannotationreach{Col.3.15}{Col}{true}\chronologyannotationreach{Col.3.16}{Col}{true}\chronologyannotationreach{Col.3.17}{Col}{true}\chronologyannotationclaim{composition.epistle-to-the-colossians}{composition}{catholic-traditional-v1}{disputed}{(Philemon; Colossians; Ephesians; Philippians), 61}{A.D. 61}; \chronologyannotationreach{Col.3.12}{Col}{true}\chronologyannotationreach{Col.3.13}{Col}{true}\chronologyannotationreach{Col.3.14}{Col}{true}\chronologyannotationreach{Col.3.15}{Col}{true}\chronologyannotationreach{Col.3.16}{Col}{true}\chronologyannotationreach{Col.3.17}{Col}{true}\chronologyannotationclaim{composition.epistle-to-the-colossians}{composition}{catholic-traditional-v1}{disputed}{some time between A.D. 61 and 63}{A.D. 61--63}\\[.25em]%
}

The passages below have several different dated subjects. Jeremias's
letter has its own relation to the deportees' departure; the received
book has a separate history of enlargement. Psalm titles identify a
tradition of singers without fixing a particular occasion. The dates
proposed for writing a Gospel concern its composition, while the scene
in which Jesus teaches concerns the narrated event. The chant's selection
of a verse does not enlarge its liturgical appointment to the complete
chapter discussed in the commentary.

% Generated. Do not edit.
% Complete orientation-v3 table; wording, dates and order belong to the corpus.
\begin{chronologysharedtable}
\chronologyorientationrow{row-a16db10bcf3bfa76}{}
\chronologyorientationrow{row-fd3920d9cd8b2ce1}{}
\midrule
\chronologyorientationrow{row-daf59e02c6c15663}{}
\chronologyorientationrow{row-d85392b5260231db}{}
\midrule
\chronologyorientationrow{row-e61fc8dcfe43e661}{}
\chronologyorientationrow{row-777243e0b53190a9}{}
\midrule
\chronologyorientationrow{row-33a9b5f385b6bb7a}{}
\chronologyorientationrow{row-42f28ec9d138025a}{}
\midrule
\chronologyorientationrow{row-ea8e2406b53ffb05}{}
\chronologyorientationrow{row-956fb69c94da31fb}{}
\midrule
\chronologyorientationrow{row-e324a33f1f79a341}{}
\chronologyorientationrow{row-5757046e665e2aca}{}
\chronologyorientationrow{row-ca2cc02c958c6d4d}{}
\midrule
\chronologyorientationrow{row-624ba2b945ad0063}{}
\chronologyorientationrow{row-ceeb5d1f4ebddba1}{}
\midrule
\chronologyorientationrow{row-b5e644ac10786f4b}{}
\end{chronologysharedtable}
{\footnotesize Psalm numbering: Vg. = Vulgate; Heb. = Hebrew.\par}


\paragraph{Setting notes.}
\noindent\textbf{Introit: Ps 84:2 (Vulgate numbering; Hebrew 85); Gradual: Psalm 43:8--9 (Vulgate numbering; Hebrew 44); Alleluia: Psalm 129:1--2 (Vulgate numbering; Hebrew 130); Offertory: Psalm 129:1--2 (Vulgate numbering; Hebrew 130)}. Israel recalls rescue and cries for mercy; individual occasions and exact places remain unidentified.

\noindent\textbf{Introit: Ps 84:2 (Vulgate numbering; Hebrew 85); Gradual: Psalm 43:8--9 (Vulgate numbering; Hebrew 44)}. Temple-singer tradition of the sons of Korah; Jerusalem is its worship horizon, with each writing site uncertain.

\noindent\textbf{Alleluia: Psalm 129:1--2 (Vulgate numbering; Hebrew 130); Offertory: Psalm 129:1--2 (Vulgate numbering; Hebrew 130)}. Israel is called to hope; the individual psalm's writing place and date are insecure. The displayed bound is Psalter-wide.

\noindent\textbf{Introit: Jer 29:11--12,14}. From Jerusalem to exiles in Babylon; Jer 29:1--3.

\noindent\textbf{Introit: Jer 29:11--12,14}. The prophet's Judah/Jerusalem ministry and later enlargement; no single final-redaction site established.

\noindent\textbf{Gospel: Matthew 13:24--30}. Jesus teaches the crowds beside the sea; the subsequent explanation is in a house (Matt 13:1--3,36).

\noindent\textbf{Gospel: Matthew 13:24--30}. The traditional original is associated with Matthew's departure from Palestine; its relation to the received Greek text is distinguished below.

\noindent\textbf{Gospel: Matthew 13:24--30}. The NABRE introduction reports a probable majority account, with Antioch in Roman Syria plausible. The dating remains a scholarly proposal.

\noindent\textbf{Communion: Mark 11:24}. Rome in the reported tradition; the instruction itself occurs within the Jerusalem/Bethania narrative.

\noindent\textbf{Communion: Mark 11:24}. Tentative anonymous-author proposal, with Syria possible; applies to Mark through 16:8, including this verse.

\noindent\textbf{Epistle: Colossians 3:12--17}. Traditional Roman captivity; recipients at Colossae and neighbouring Laodicea. The modern introduction leaves the prison site and authorship uncertain.


\subsection{Psalms: attribution, occasion and collection}
\label{hist-psalms}
Psalms 43 and 84 have the sons-of-Korah attribution in the received
Psalter. Their words concern communal rescue and the Lord's land,
but the headings identify no particular event that fixes either poem's
occasion. Psalm 43 recalls the ancestors' victories and then laments
later defeat; the recalled victories cannot simply date the lament.
Psalm 84's thanksgiving for restoration is followed by petitions for
further restoration. Psalm 129 is a gradual canticle without a named
author or composition site.

The historical \work{Catholic Encyclopedia}'s account of the Korahite
psalms supplies the composition horizon printed above. Attribution to
a guild permits the work of collection as well as composition by a
single author. The temple-singer tradition gives Jerusalem as a worship
horizon, while leaving the actual writing place of each poem uncertain.
The USCCB Psalms introduction stresses the difficulty of secure dating
for individual psalms. Its latest collection-wide boundary, also printed
above, gives a limit rather than an individual occasion.\footnote{Walter
Drum, \href{https://www.newadvent.org/cathen/12533a.htm}{``Psalms,''}
\work{Catholic Encyclopedia} XII (1911), titles/authorship, sons-of-Korah
paragraph; USCCB, \href{https://bible.usccb.org/bible/psalms/0}{Psalms
introduction}, ``Historical Settings.''}

Augustine's accounts of sin's captivity, the martyrs' trials and the
resurrection are theological expositions of these psalms. They do not
identify the poems' composition dates. His NPNF headings use the
Hebrew psalm number while the present study follows the Missal's
Vulgate numbering. The particular English Bible used here numbers the
Gradual's verses as 43:8--9. Conversion to an English Bible under
Psalm 44 requires attention to that edition's verse numbering as well
as to its psalm number.

\subsection{Jeremias: the letter and the enlarged book}
\label{hist-jeremias}
Jeremias 29:1--3 identifies the letter's addressees and route: from
Jerusalem to people carried to Babylon after Jechonias's departure.
That relationship, printed in the table, dates the letter relatively.
The letter's first unit ends at verse 23. The Introit selects verses
11, 12 and 14 rather than appointing the entire unit.

Faulhaber's account distinguishes the book's first authentic recension
from its later enlargement. The original recension is discussed in
relation to Jeremias 36; it is a different subject from the later
constituent letter to the exiles. Received-book formation is consequently
printed as a separate relation. No single final-redaction place is
established by this account.\footnote{Michael Faulhaber,
\href{https://www.newadvent.org/cathen/08334a.htm}{``Jeremias,''}
\work{Catholic Encyclopedia} VIII (1910), ``Analysis of contents'' and
``Literary criticism of the book''; Jer 29:1--3, 36.}
The Introit's psalm constituent has its own attribution and chronological
questions. A common chant cannot give these separately composed texts
one composition date.

\subsection{Matthew: teaching by the sea and the written Gospel}
\label{hist-matthew}
Matthew 13:1--3 locates Jesus with crowds beside the sea, and verse 36
places the subsequent explanation in the house. The teaching occasion
has no resolved calendar date in the evidence above. The sowing and
harvest depicted in the parable are its narrated comparison rather
than dated events in an identified farmer's life.

Jacquier reports several traditional estimates for Matthew's composition.
His account distinguishes an original Hebrew or Aramaic form from
the received Greek Gospel. The \work{Catholic Encyclopedia}'s ``New
Testament'' article supplies another estimate for the original and
leaves the date of its rendering into Greek indefinite. The early
and late alternatives above remain distinct proposals rather than
one continuous range of confidence.\footnote{E. Jacquier,
\href{https://www.newadvent.org/cathen/10057a.htm}{``Gospel of St.
Matthew,''} \work{Catholic Encyclopedia} X (1911), ``Date and place
of composition'' and concluding Biblical Commission account;
\href{https://www.newadvent.org/cathen/14530a.htm}{``The New Testament,''}
XIV (1912), origin of the books, Matthew paragraph.}

The modern Catholic-critical comparison concerns the received Greek
Gospel. The USCCB introduction describes a probable majority account
of an unnamed Greek evangelist using Mark and other traditions. It
regards a later decade as probable and Antioch as a plausible place,
while the generated table states the open lower bound retained for
this comparison. The probabilistic authorship, date and place proposal
concerns this Greek composition; it neither dates an Aramaic predecessor
nor assigns a year to Jesus's teaching.\footnote{USCCB,
\href{https://bible.usccb.org/bible/matthew/0}{Matthew introduction},
final authorship/source/date and place paragraphs beginning ``The questions
of authorship'' and ``As for the place.''}

\subsection{Mark: reported tradition and a tentative modern proposal}
\label{hist-mark}
Mark 11:11--12 and 20--27 locate the instruction within journeys between
Jerusalem and Bethania. That narrative location differs from the place
where the Gospel was written. MacRory reports the Roman composition
tradition and several historical dating proposals, explaining his own
conclusion separately. The ``New Testament'' article reports another
range. The Biblical Commission's response adds a relative order and
bound, retained above in its own terms.\footnote{Joseph MacRory,
\href{https://www.newadvent.org/cathen/09674b.htm}{``Gospel of St.
Mark,''} \work{Catholic Encyclopedia} IX (1910), place/date/destination
paragraphs; ``The New Testament,'' XIV (1912), Mark paragraph;
Pontifical Biblical Commission, \work{De auctore, de tempore compositionis,
et de historica veritate Evangeliorum secundum Marcum et secundum Lucam},
AAS 4 (1912), pp.~463--465, especially responses V--VI on p.~464.}

The USCCB introduction distinguishes the John Mark/Petrine tradition
from a tentative modern account of an anonymous Hellenistic Jewish
Christian author, possibly writing in Syria. Its dating proposal remains
open and expressly tentative. It applies to the Gospel through Mark
16:8, within which the Communion verse falls; the added endings are
a separate textual subject. Neither the traditional alternatives nor
the modern comparison dates the fig-tree episode itself.\footnote{USCCB,
\href{https://bible.usccb.org/bible/mark/0}{Mark introduction}, final
historical paragraphs, complete inspected text extraction, lines 72--75.}

\subsection{Colossians: the captive writer and the recipients}
\label{hist-colossians}
The traditional orientation places Paul in Roman captivity and the
recipients at Colossae, with a neighbouring community at Laodicea.
The letter names its recipients at the beginning and names messengers
and neighbouring churches in its conclusion (Col 1:2; 4:7--18).
The writer's captivity and the recipients' location answer different
geographical questions. Prat and Aherne supply the two traditional
composition alternatives printed above.\footnote{Ferdinand Prat,
\href{https://www.newadvent.org/cathen/11567b.htm}{``St. Paul,''}
\work{Catholic Encyclopedia} XI (1911), chronology summary and Captivity
Epistles; Cornelius Aherne,
\href{https://www.newadvent.org/cathen/04131b.htm}{``Epistle to the
Colossians,''} IV (1908), opening captivity paragraph and recipients.}

The USCCB introduction leaves both prison location and precise composition
date uncertain. It includes Roman custody among proposals, with Caesarea
and Ephesus also considered, and reports the possibility of a pupil
writing in Paul's name. That alternative receives no numeric date here.
The table's Roman description is attached to the traditional account,
while the uncertainty qualifies the wider historical question.\footnote{USCCB,
\href{https://bible.usccb.org/bible/colossians/0}{Colossians introduction},
final historical paragraph.}
